\cotefile{93}
\scanpage{47}
peut énoncer qu'il y a un isomorphisme canonique (fonctoriel, compatible
aux extensions de la base)
\[
\deleted{\unsure{\underline E(}}\handwritten{\mathbb D}^{*}(A)=\underline{\mathrm{Lie}}(\underline E(A^{*}))\ ,
\]
la structure d'extension de $\mathbb D^{*}(A)$ n'étant autre que la structure\inserted{\handwritten{sur \uline{Lie}}}
déduite de la structure d'extension de $\underline E(A)_S=E(A)$.

Ceci posé, revenant aux conditions du théo\deleted{\unsure{èr|rè}}rème de déformations, on
voit, par un sorite sur l'exponentielle sur lequel nous reviendrons, que
la donnée d'un prolongement $\mathrm{Fil}^{1}\handwritten{\subset}\mathbb D^{*}(A)_{S'}=\underline{\mathrm{Lie}}(\underline E(A^{*})_{S'})$ localement facteur
direct prolongeant $\underline{\mathrm{Lie}}(\underline{\underline{t}}^{\handwritten{\vee}}_A)=\underline t^{\handwritten{\vee}}_A\handwritten{\subset}\mathbb D^{*}(A)_S=\underline{\mathrm{Lie}}(\underline E(A^{*})_{S'}\handwritten{\mid} S)$ équivaut à
la donnée d'un sous-groupe \inserted{lisse} a structure vectorielle de $\underline E(A^{*})_{S'}$ qui prolonge le sous-groupe $\underline t_A^{\handwritten{\vee}}$ de $\underline E(A^{*})_S$. Si $L$ est ce sous-groupe, alors il est
immédiat que \deleted{$\unsure{A^{*}(|E^{*}(}$}
\[
\underline E(A^{*})_{\deleted{\unsure{t}}\handwritten{S'}}/L
\]
est un schéma abélien sur $S'$, dont le schéma abélien dual est le $A'$ cherché. (NB Il serait plus naturel ici de travailler avec \deleted{$\unsure{D_{*}(A)=D(A)|\mathbb D^{*}(A)=\mathbb D(A)}$}
\[
\mathbb D_{*}(A)=\underline{\mathrm{Lie}}\ \underline E(A)\ \overset{\text{isom can}}{\handwritten{\xrightarrow{\sim}}}\ \mathbb D^{*}(A^{*})^{\handwritten{\vee}}
\]
plutôt qu'avec $D^{*}(A)$ pour énoncer le théorème de déformations; l'équivalence des deux points de vue provient de l'accouplement parfait
\[
\mathbb D^{*}(A)\otimes\mathbb D^{*}(A^{*})\handwritten{\longrightarrow}\underline O_{S_{\mathrm{c}\deleted{\unsure{\mathrm{ii}}}\mathrm{risnilp}}}
\]
compatible avec les filtrations.\handwritten{)}

\uline{Remarque}. En t\deleted{\unsure{éh}}héorie des groupes de BT, nous donnerons deux constructions de $\handwritten{\mathbb D}^{\handwritten{*}}(G)$, inspirées des deux constructions indiquées dans le cas des
schémas abéliens: l'une "cohomologique", l'autre via une "extension vectorielle universelle" $\underline E(G)$. La construction de $\underline E(G)$ peut aussi s'obtenir
par deux méthodes, l'une cohomologique, l'autre"par l'exponentielle" étant
directement adaptée à la démonstration du théorème de déformations.
% archive sheet number: 46

\scanpage{48}
\handwritten{3}) Relations entre les deux théories (pour schémas abéliens et pour
\deleted{\unsure{c)|2)} Relations} groupes de BT). Supposons que $p$ soit localeme\deleted{\unsure{j}}nt nilpotent sur $S$. Travaillons \deleted{On suppose} à nouveau avec le site cristallin de Berthel\deleted{l}ot (pas le
nilpotent). Le résul\deleted{\unsure{lt}}tat essentiel sera, pour un s\deleted{\unsure{v}}chéma abélien variable\deleted{\unsure{,}}
$A$ sur $S$, un isomorphisme fonctoriel compatible aux changements de base
\[
\underline D^{\handwritten{*}}(A(\infty))\handwritten{\simeq} R^{1}f_{\mathrm{cris}\,\handwritten{*}}(\underline O_{A_{\mathrm{cris}}})\quad ,
\]
cet isomorphisme induisant, pour les valeurs des deux membres sur $S$, un
isomorphisme compatible aux filtrations \qquad et \qquad (compte tenu des
isomorphismes $\underline t_{A(\infty)}\handwritten{\simeq}\underline t_A$ et idoine pour $A^{*}$).\inserted{\handwritten{-- C'est en fait en postulant
un tel isomorphisme qu'on arrive par voie heuristique à une définition de $\mathbb D^{*}(G)$ pour un groupe de BT $G$ quelconque}}
% slanted: the insertion continues in the left margin after « définition ».

Ceci impliquera en particulier le fait que, pour un \deleted{voisinage} épaississement à p. div. loc. nilp. \deleted{de} $S'$ de $S$, il y a identité de la théorie
des prolongements infinitésimaux de $A$, et de $A(\infty)$. De eeci, par dévissage
sur l'ordre de nilpot\deleted{\unsure{t}}ence, on peut déduire par exemple immédi\deleted{\unsure{q}}atement le
théorème de Serre-Tate déjà annoncé plus haut\handwritten{:}\inserted{\handwritten{ méthode suivie par
Messing ds sa thèse.}}

\handwritten{[}On peut aussi définir une extension vectorielle universe\deleted{\unsure{e}}lle \deleted{$\underline E$} $E(G)$
d'un groupe de BT, et on trouvera une isomorphisme canonique etc
\[
E(A(\infty))\handwritten{\simeq} E(A)(\infty)
\]
donnant naissance, par le procédé habituel, à un isomorphisme plus
général de cristaux en \deleted{G}groupes
\[
\underline E(A\handwritten{(}\infty\handwritten{))}\handwritten{\simeq}\underline E(A)(\infty)\quad\handwritten{....}\qquad\handwritten{]}
\]

\deleted{4) Passage à la limite}

\deleted{\handwritten{4}) Problème\unsure{.|:} trouver le foncteur des F-cristaux filtrés vers \deleted{\unsure{modules}}
faisceaux $p$-adiques.}
% archive sheet number: 47

\scanpage{49}
Rabiots pour l'exposé précédent:

a) \deleted{\unsure{Définition}}\unsure{Déscription|Description} de la filtration canonique d'un F-V-cristal dans
le cas d'une base \inserted{un anneau \uline{parfait}} \handwritten{(}de car. $p$\handwritten{)}: elle est définie par un sous-module
$\mathrm{Fil}^{1}\ M_{\unsure{o|0}}$ \deleted{($M$ considéré comme module localement libre de type fini sur
$W=\underline W(A)$)} tel que l'on ait
\[
\mathrm{Fil}^{1\ (p)}=\mathrm{Ker}\ F_M=\mathrm{Im}\ V_M\quad ,
\]
ce qui le définit de façon unique grâce à l'hypothèse sur $A$ (NB tout
\deleted{\unsure{c|M}}Cristal de Dieudonné est alors "admissible").

b) Iso\deleted{cristaux}groupes \inserted{de BT} (devrait faire un sous-paragraphe). Définition
sur une base quelconque: on localise \deleted{pour $Z$} en remplaçant
$\mathrm{Hom}\ (G,H)$ par $\mathrm{Hom}(G,H)\otimes_{\underline Z_p}\underline Q_p=\mathrm{Hom}\otimes_{\underline Z}\underline Q$. (Définition provisoire; elle
\deleted{\uline{Proposition}} a l'inconvénient de ne pas fournir un champ, et devrait
peut-être être reconsidérée ...). On définit de même les isocristaux
de Dieudonn\deleted{\unsure{è}}é, et on obtient donc un foncteur
\[
\mathrm{IsoBT}(S)\qquad\qquad\mathrm{Iso}\deleted{\mathrm{cris}}\text{F-V-cris}(S)
\]
($p$ localement nilpotent sur $\deleted{\unsure{J}}S$). On trouve

\uline{Proposition} (prouvée plus tard): \deleted{Supposons $S$ quasi-compact, et soit
soit \unsure{$S'$|$S$|$S_0$} un \unsure{nilidéal d'id|idéal d'id}} \inserted{Soit $S\qquad S'$ une immersion nilpoten}te d'Idéal $J$ annulé par une puissance de
$J$, alors le foncteur restriction
\[
\mathrm{IsoBT}(S')\handwritten{\longrightarrow}\mathrm{IsoBT}(S)
\]
est une équivalence de catégories.

On en conclut l'énoncé VII 10.8.2 de mes notes de Pise.\inserted{\handwritten{Expliciter le cas d'un anneau de val. discrète d'inégales car.}}
\margin{\handwritten{\uline{\unsure{Th.|Th}}}}
On peut se poser le \deleted{question de l'image essentielle}

\uline{Problème} 4. Quel est l'image essentielle dudit foncteur ? Suppposant
$A$ noethérien séparé et complet pour la topologique $p$-adique pour fixer
les idées, y a-t-il un critère valuatif d'appartenance à l'image essentielle ? Si $A$ est un anneau de val. discrète \deleted{\unsure{(}}d'inégales caractéristiques, l'image essentielle est elle un ouvert de \deleted{\unsure{z}}Zariski de la grassma-
% archive sheet number: 48

\scanpage{50}
\mspage{- 36 -}
n\textsuperscript{o} 393

\deleted{que $\|u_n\|<C$, $x_n=\displaystyle\sum_{k=1}^{n}\gamma(1-\gamma)^{k}u_k+(1-\gamma)^{n}e$ admette un inverse $t_n$, et
$\|t_n z-t_{n-1}z\|<\dfrac{\delta}{2^n}$, $\|t_1 z-z\|<\dfrac{\delta}{2}$. Posant $y_n=t_n z$, les $y_n$ tendent vers
une limite $y\in A$, $x_n$ tend vers $x=\displaystyle\sum_{k=1}^{\infty}\gamma(1-\gamma)^{k-1}u_k$ ; $x$ et $y$ ont les
propriétés requises).}

\deleted{23) a) Soient $E$ un espace de Banach, $f:\handwritten{\mathbb C}\longrightarrow E$ une fonction entière, et
$\vartheta\longrightarrow p(\vartheta)=\displaystyle\sum_{\mu=k}^{m}c_{\mu}e^{\mu i\vartheta}$ un polynôme trigonométrique, ($c_k\ne0$). On suppose
que $|p(\vartheta)|\ \|f(re^{i\vartheta})\|\leq Mr^{\alpha}$ pour $r\geq r_{\unsure{o|0}}$, ($M$, $r_o$, $\alpha$ étant des constantes $\geq0$).
Alors $f$ est un polynôme de degré $\leq\alpha$. (Soient $f(\zeta)=\displaystyle\sum_{\nu=0}^{\infty}a_{\nu}\zeta^{\nu}$, et
$J(r,p)=r^{-p}\displaystyle\int_0^{2\pi}p(\vartheta)f(re^{i\vartheta})e^{-(k+p)i\vartheta}\,d\vartheta$. On a $J(r,p)\longrightarrow0$ quand $r\longrightarrow+\infty$
si $p>\alpha$, et d'autre part $\lim J(r;p)=c_k\,a_p$).}

\deleted{b) Prouver de même que, si $\alpha$ est entier et si $r^{-\alpha}|p(\vartheta)|\ \|f(re^{i\vartheta})\|\longrightarrow0$
quand $r\longrightarrow+\infty$, $f$ est un polynôme de degré $<\alpha$.}

\deleted{c) Soient $A$ une algèbre de Banach unipollente, $x$ un élément quasi-nilpotent
de $A$. On suppose que, pour un certain entier $n>0$, on a $r^{-n}|\cos^{n}\vartheta|\ \|(e-\lambda x)^{-1}\|\longrightarrow
\longrightarrow0$ quand $r\longrightarrow+\infty$, ($\lambda=re^{i\vartheta}$). Alors $x^n=0$. (Utiliser b)).}

\margin{\handwritten{\P}}
\deleted{24) a) Soient $E$ un espace de Banach, $x:\handwritten{\mathbb C}-\{1\}\longrightarrow E$ une fonction entière
de $1/(\zeta-1)$, $x(\zeta)=\displaystyle\sum_{n=0}^{\infty}a_n\,\zeta^n$, $x(\zeta)=\displaystyle\sum_{n=0}^{\infty}b_n\,\zeta^{-n}$ les développements de $x$
pour $|\zeta|<1$, $|\zeta|>1$ respectivement. On suppose $\|a_n\|=\unsure{o|\mathcal O}(n^\alpha)$,
$\|b_n\|=o(n^\alpha)$, ($\alpha>0$). Alors $x(\zeta)$ est un polynôme en $1/(\zeta-1)$ de degré
$<\alpha+1$. (Soit $\varepsilon>0$. On a $\|x(\zeta)\|\leq\varepsilon(1-r)^{-1-\alpha}$ pour $r_{\varepsilon}\leq r=|\zeta|<1$,
$\|x(\zeta)\|\leq\varepsilon(1-r^{-1})^{-1-\alpha}$ pour $1<r<r'_{\varepsilon}$. Posant $1-\zeta=(\omega+\frac12)^{-1}$, $\omega=Re^{i\vartheta}$,
montrer que, pour $0\leq r<1$ et $R\geq1$, on a $(1-r)^{-1}\leq\dfrac94\,\dfrac R{\cos\vartheta}$ et que, pour
$1<r$ et $R\geq1$, $(1-r^{-1})^{-1}\leq\dfrac94\,\dfrac R{\cos\vartheta}$. Poser $y(\omega)=x(\zeta)$ et montrer que}
% archive sheet number: 49

\scanpage{51}
nienne sur le corps des fractions \deleted{de} $L$ de $K$ ?

On notera que même dans le cas où $A=W(k)$, $k$ alg. clos, le
foncteur en question n'est pas \deleted{essentiellem} essentiellement surjectif.

Notons le

\uline{Problème} 5. \inserted{($A$ anneau de val. discrète complet d'inégales car., $k$ parf)} Expliciter le foncteur naturel qui va de la sous-catégorie pleine précédente de la catégorie des couples \deleted{d} $(M,\mathrm{Fil})$ d'un module
de Dieudonné $M$ sur \deleted{$W$} $k$ (i.e. \deleted{sur} module sur $W$) et d'une filtration à
deux crans de $M\otimes_W L$, vers les \deleted{mo}dules galoisiens sur $L$ (représent\unsure{a}tions
\inserted{vec\deleted{\unsure{toir}}torielles des dimension finie sur $\underline Q_p$ de}
de $\mathrm{Gal}(\bar L/L)$)\deleted{dans des \unsure{représen|espaces} modules libres de type fini} Peut-on même
définir ce foncteur \deleted{sur la sous-catégorie} de façon naturelle sans se
borner à cette sous-catégorie pleine,\deleted{\unsure{a}}en imposant seulement des conditions de rangs sur le gradué, mettons ?

\uline{Problème} 6. Revenant à un $A$ comme dans Pb 4, peut-on trouver un foncteur \deleted{d'une \unsure{c}} de la \deleted{\unsure{a}}catégorie des couples $(M,\mathrm{Fil}^{1})$ d'un isocristal
de Dieudonné \inserted{admissible} sur $A_{\unsure{o|0}}=A/pA$ (ou encore, sur $A_{o\ \mathrm{red}}$ \deleted{\illegible{1}}) et d'une filtration en deux crans par modules localement libres de $(M_A)_p$, ayant le
rang qu'il faut\deleted{,} (et sou\deleted{s}mis éventuellement à d'autres conditions
ouvertes \inserted{au sens Zar}, assurant que le couple envisagé provient d'un isogroupe de
BT sur $A$),vers les faisceaux $p$-adiques constants tordus sur $U=\mathrm{Spec}(A_p)$,
qui exprime le foncteur du théorè\deleted{\unsure{p}}me (VII 10.8.2) précédent.

\deleted{4) Question des relations entre cohomologie cristalline et $p$-adique.}

\deleted{\deleted{Problème} \uline{Conjecture} 7. On peut, pour tout anneau de valuation discrète complet $V$ d'inégales caractéristiques, de corps \deleted{des fractions} ris\deleted{\unsure{u}}\unsure{i|u}duel parfait (pour simplifier), de corps des fractions $L$, trouver un
foncteur \deleted{\unsure{san}}canoniques sur une sous-catégorie pleine convenable des couples $(M,\mathrm{Fil}^{\deleted{1}})$ \deleted{envisagée} dans le Pb 5 vers les}
% archive sheet number: 50

\scanpage{52}
\mspage{- 35 -}
n\textsuperscript{o} 393

\deleted{21) Soient $A$ une algèbre de Banach unipollente, $\Delta$ une partie ouverte de $\handwritten{\mathbb C}$,
$\lambda\longrightarrow R(\lambda)$ une application de $\Delta$ dans $A$ telle que :
\[
R(\lambda)-R(\mu)=-(\lambda-\mu)R(\lambda)R(\mu)\tag*{(1)}
\]
quels que soient $\lambda$, $\mu\in\Delta$.}

\deleted{a) Soit $\lambda_{\unsure{o|0}}\in\Delta$. Pour tout $\lambda\in\Delta$ tel que $|\lambda-\lambda_o|\|R(\lambda)\|<1$,
on a :
\[
R(\lambda)=\sum_{n=0}^{\infty}(\lambda_o-\lambda)^n\,R(\lambda_o)^{n+1}.
\]}

\deleted{b) Pour tout $\lambda\in\Delta$, $\dfrac{d^n}{d\lambda^n}R(\lambda)=(-1)^n\,n!\,R(\lambda)^{n+1}$.}

\deleted{c) Si $R$ est holomorphe à l'infini, il existe $z$, $j$, $x\in A$ tels que $z^2=0$,
$j^2=j$, $zj=jz=0$, $x\in jAj$, et $R(\lambda)=z+jR(\lambda,x)$ pour $|\lambda|$ assez grand.
(Ecrire $R(\lambda)=\displaystyle\sum_{n=0}^{\infty}c_n\,\lambda^{-n}$ pour $|\lambda|>\lambda_o$, et exprimer que $R$ satisfait à (1).
On peut poser $c_o=z$, $c_1=j$, $c_2=x$).}

\deleted{d) Pour que $R$ soit la résolvante d'un élément de $A$, il faut et il suffit que
$R(\lambda_o)$ soit inversible pour un $\lambda_o\in\Delta$.}

\deleted{e) Pour tout $x\in A$, la fonction $\lambda\longrightarrow\displaystyle\sum_{n=0}^{\infty}(\lambda_o-\lambda)^n\,x^{n+1}$ satisfait à (1)
pour $|\lambda-\lambda_o|\ \|x\|<1$.}

\margin{\handwritten{\P}}
\deleted{22) a) Soient $B$ une algèbre de Banach unipollente, $w\in B$, $C>1$ et $\varepsilon>0$. Il
existe $\eta>0$ tel que les conditions $u\in B$, $\|u\|<C$, $\|uw-w\|<\eta$,
$0\leq\gamma\leq1/4C$ impliquent $\|yw-u\|<\varepsilon$, où $y=[(1-\gamma)e+\gamma u]^{-1}$.}

\deleted{b) Soit $A$ une algèbre de Banach. On suppose qu'il existe une constante $C$ telle
que, pour tout $\varepsilon>0$ et tous $x_1,\ldots,x_n\in A$, il existe $u\in A$ avec $\|u\|<C$,
$\|ux_1-x_1\|<\varepsilon,\ldots,\|ux_n-x_n\|<\varepsilon$. Pour tout $z\in A$, il existe $x$, $y\in A$
tels que $z=xy$, $y\in\overline{Az}$, $\|z-y\|<\delta$. (On peut supposer $C>1$. Soit $\gamma=1/4C$
Appliquant a) dans $\widetilde A$, déterminer par récurrence une suite $u_1$, $u_2$,\ldots\ dans $A$ telle}
% archive sheet number: 51

\scanpage{53}
\margin{\handwritten{\unsure{suite de|mieux de} pentes de Hodge}}
% slanted: leader runs from the marginal words to the first primed lambda.
\[
\begin{aligned}
&\handwritten{\left\{\begin{aligned}
\sum_{1}^{i}\lambda'_j&\leq\deleted{\unsure{\vee}}\sum_{1}^{i}\lambda_j\\
\sum_{\unsure{n|N|1}}^{\deleted{\unsure{n}}}\lambda'_j&=\sum\lambda_j
\end{aligned}\right.}\\[-1pt]
&\qquad\handwritten{\overset{\unsure{''|=}}{\sum}\,i h_i^{\prime}}
\end{aligned}
\]

\noindent\rule{\linewidth}{0.3pt}\par

\deleted{\handwritten{\uline{\unsure{Sur un th. de Tate.}}}} Sur un théorème de Tate. (Devrait venir avec foncteur mystérieux !)

C'est le théorème suivant (Driebergen ....):

\uline{Th.} (Tate): Soit $S$ un schéma noethérien normal intègre dont le corps des
fonctions \inserted{$k(\handwritten{\eta})$} est de car. zero, Alors le foncteur
\[
\mathrm{BT}(S)\handwritten{\longrightarrow}\mathrm{BT}(\handwritten{\eta})
\]
est pleinement fidèle.

Ainsi, un groupe de BT sur $S$ est c\deleted{\unsure{b}}onnu \inserted{à isom unique près} quand on connait la représentation $p$-adique qu'\deleted{elle}il définit de $\handwritten{\pi}_{1}(\handwritten{\eta})=\mathrm{Gal}(\overline{k(\handwritten{\eta})}/k(\handwritten{\eta}))$. Prenant d'abord
le cas où $S$ est le spectre d'un anneau de valuation discrète complet $V$
d'inégales car. , de corps \deleted{des} résiduel $k$ parfait, la question se pose donc

a) De caractériser les modules galoisiens qu'on peut obtenir (Tate donne
des conditions nécessaires: les modules en question doivent être du type de
"Hodge-Tate" );

b) De trouver un foncteur "explicite" sur la catégorie de ces modules galoisiens, vers les F-cristaux sur $k\quad M$ muni d'une filtration de $M\otimes_W\deleted{\unsure{V}}L$.

\handwritten{\uline{Problème}} Le théorème de Tate suggère de plus que le "foncteur mystérieux"
qui va des F-cristaux sur $k$ \deleted{\unsure{m}}filtrés sur $L$, vers les $\underline Q_p$-modules galoisiens
sur $L$, est pleine\deleted{\unsure{n}}ment fidèle, et qu'on doit pouvoir donc \deleted{définir un foncteu}
expliciter un foncteur quasi-\deleted{\unsure{c}}inverse (non moins mystérieux), allant d'une
sous-catégorie pleine convenable des $\underline Q_p$-modules \deleted{\unsure{G}}galoisiens sur $L$ (sans
doute contenue dans celle des modules galoisiens du type de Hodge-Tate)
vers les F-cristaux sur $k$ filtrés sur $L$.

D'autre part, on conjec\deleted{\unsure{y}}ture que dans le th. de Tate ci-dessus, la
restriction au cas où \deleted{$S$}\inserted{$k(\handwritten{\eta})\handwritten{=L}$} est de car. nulle est inutile, Le seul cas qui offre
% archive sheet number: 52

\scanpage{54}
un Pb. est celui où $k(\handwritten{\eta})=L$ est de car.$p$ (cas d'égales caractéristiques).
On peut alors penser à attaquer le Pb via la théorie de Dieudonné (supposan\deleted{\unsure{s}}t résolu le Pb 1) et via le

\uline{Problème} n\textsuperscript{\uline{o}} ... Soit $S$ un schéma noethérien normal \deleted{de car} intègre d\deleted{\unsure{a}}e car.
$p$, de \deleted{\unsure{o}}point générique $\handwritten{\eta}$. Est-il vrai que le foncteur de restriction
\[
\text{Cris loc lib }(S)\handwritten{\longrightarrow}\text{ Cris loc lib }(\handwritten{\eta})
\]
(ou du moins le foncteur analogue pour F-cristaux) est pleinement fidèle ?

On est ramené facilement au cas où $S$ est le spectre d'un anneau de
\deleted{\unsure{g}}la forme $k\handwritten{[[}t\handwritten{]]}$, avec $k$ alg. clos, auquel cas l'énoncé s'explicite en
termes de Modules avec connexions formelles .... Il semble que si la réponse au Pb précédent est affirmative, \deleted{ainsi que} et si on sait faire la
théorie de Dieudonné\unsure{n} sur un anneau \uline{parfait} (avec \uline{équivalence} de catégories
\deleted{\unsure{q}}alors on aura démonstration du th. de Tate en égales car.
% archive sheet number: 53

\scanpage{55}
$\chi$ et $\chi_!$
% archive sheet number: 54

\scanpage{56}
$k$ corps
\diagram{\small\(\begin{array}{l l l l}
\unsure{\mathcal E|\mathscr E}_{\chi_!}&\text{espaces topologiques \inserted{loc. compacts} tels que}&H^*_{!}(\deleted{\unsure{k}}X,k)&\text{de dim finie sur }k\\
&&\multicolumn{2}{l}{\quad\text{coh. \deleted{\unsure{faisceautique|faisceaux}} \inserted{singulière}}}\\
\mathcal E_{\chi}&\rule{3.8cm}{0.3pt}&H^*(X,k)&\rule{2.4cm}{0.3pt}\\
&&\multicolumn{2}{l}{\quad\text{coh. singulière}}
\end{array}\)}

\textcircled{a} \deleted{\unsure{$X$|Si $X$} \unsure{est} \unsure{cpt}} Si $X$ compact, $X\in\mathcal E_{\chi_!}\Longleftrightarrow X\in\mathcal E_{\chi}$, et alors
\[
\chi_!(X)=\chi(X).\qquad\qquad\text{trivial.}
\]

\textcircled{b} Si $Y$ fermé dans $X$ \inserted{loc. cpt}, alors si deux des espaces $X,Y,X-Y$
sont dans $\mathcal E_{\chi_!}$, le troisième aussi, et on a
\[
\chi_!(X)=\chi_!(Y)+\chi_!(X-Y)
\]
(résulte de suite exacte de cohomologie à supp. cpts)

\textcircled{c} Si $X',X''$ \deleted{\unsure{de}} fermés dans $X$ de réunion $X$, alors si trois
parmi les quatre espaces $X,X',X'',X'\cap X''$ sont dans
$\mathcal E_{\chi}$, le quatrième aussi, et on a
\[
\chi(X)=\chi(X')+\chi(X'')-\chi(X'\cap X'')
\]
(résulte de suite exacte de Mayer-Vietoris

\deleted{\textcircled{d}} [Même $X,X'\cap X''\in\mathcal E_{\chi}\Longrightarrow X'\unsure{;|,}X''\in\mathcal E_{\chi}$]

\textcircled{d} Si $X$ \deleted{\unsure{est} \illegible{1} $X$}\inserted{est dominé} homotopiquement par $X'$
\diagram{(\begin{tikzcd}[ampersand replacement=\&,column sep=large]
X\arrow[r,shift left=0.6ex,"i"]\&X'\arrow[l,shift left=0.6ex,"j"]
\end{tikzcd}\quad $ji\unsure{\simeq|\sim}\mathrm{id}_{X}$)}
alors $X'\in\mathcal E_{\chi}\Longrightarrow\deleted{\unsure{X'}}X\in\mathcal E_{\chi}$, Donc
$(X\sim X')\unsure{\Longleftrightarrow|\Longrightarrow}(X\in\mathcal E_{\chi}\Longleftrightarrow X'\in\mathcal E_{\chi})$, et de plus
on a alors
\[
\chi(X)=\chi(X')
\]

\textcircled{e} \deleted{\unsure{Si} $X$ est \unsure{connexe} ; un pt,} $\mathrm{Card}\,X\leq1\Longrightarrow X\in\mathcal E_{\chi}$, de
plus on a
\[
\chi(\mathrm{pt})=1
\]
(de plus $\chi(\varnothing)=0$, \unsure{grâce à|qui suit} c) p.ex) \qquad (trivial)
% archive sheet number: 55

\scanpage{57}
\mspage{10.}
5.2. \uline{Evaluation de la politique de recherche des établissements}

La D.E.S.U.R. ne sera crédible auprès de la D.G.R.S.T.
et par conséquent n'obtiendra des crédits suffisants que si elle est capable de
porter un jugement sur la valeur des résultats de recherche de chaque établissement
et sur la façon dont les établissements conduisent leur politique de recherche.
Actuellement les services de la D.E.S.U.R. disposent de nombreuses données sur les
établissements :

\begin{itemize}[label={-}]
\item le nombre des thèses soutenues
\item le nombre et la qualité des publications
\item le nombre d'associations avec les autres établissements de recherche
\item le nombre des contrats
\item évolution de la carrière du personnel
\end{itemize}

Le danger est grand de voir s'instituer un système purement
numérique, alimenté par ces critères, et fournissant, par un ordinateur, le\unsure{:|;}montant
des crédits de l'université. Juger la recherche conduite par un établissement est
une opération délicate ; il est donc envisagé de placer, auprès du directeur des
enseignements supérieurs et de la recherche un groupe de personnalités indépendantes
capables de faire la synthèse des renseignements recueillis sur un établissement
pour porter un jugement global qui restera confidentiel.

5.3. \uline{La responsabilité de la D.E.S.U.R.}

Les deux branches du Secrétariat d'Etat qui supportent la
recherche dans les établissements d'enseignement supérieur, la D.E.S.U.R. et le
C.N.R.S. doivent pratiquer une politique différente , quoique étroitement complémentaire. Le C.N.R.S. est une entreprise de recherche centralisée, qui agit par
discipline ; il aide, dans un tel établissement tel laboratoire parce que cela est
utile au programme de recherche. La D.E.S.U.R. agit par établissement ; elle aide
chaque établissement à développer le programme qu'il propose, sous réserve qu'il
contribue à l'une des missions spécifiques de la recherche universitaire.

Cela suppose la mise au point d'une politique de répartition
des moyens complémentaires (\S\ 4.2.) qui en s'ajoutant aux moyens de base (\S4.1.)
conditionnent la politique de recherche de l'établissement.

\hfill .../...
% archive sheet number: 56

\scanpage{58}

\uline{Conclusions premières} (on pose $\overline{\mathbb R}=[-\infty,+\infty]$

\margin{\textcircled{1} Soit $X',X''\in\unsure{\mathcal E|\mathscr E}_\chi$ (resp. $\in\mathcal E_{\chi_!}$) $\Leftrightarrow X'\sqcup X''\in\mathcal E_\chi$ (resp. $\in\mathcal E_{\chi_!}$), \deleted{résulte de c et \unsure{e}} et alors $\chi(X)=\chi(X')+\chi(X'')$ (resp. $\ldots$) résulte de c) et e) ($\varnothing\in\mathcal E_\chi$) resp. de b)}

\textcircled{\unsure{1|2}} \deleted{\unsure{Soit|Pour} \unsure{$I$}} \inserted{$\forall n\in\mathbb N$}, $R^n,\overline R^n,B^n,S^n\in\mathcal E_\chi\cap\mathcal E_{\chi_!}$, de plus
\[
\left\{\begin{aligned}
\chi(\mathbb R^n)=\chi(\overline R^n)&=\chi(B^n)=1 & \chi(S^n)&=\chi_!(S^n)=1+(-1)^n\\
\overset{\Vert}{\chi_!}(\overline R^n)&=\overset{\Vert}{\chi_!}(B^n)\\
\chi_!(R^n)&=(-1)^n
\end{aligned}\right.
\]

\uline{Dém} $R^n,\overline R^n,B^n\sim\mathrm{pt}$, donc par d) et e) ils sont $\in\mathcal E_\chi$ ; comme $\overline R^n,B^n$ compacts, par b) ils sont aussi $\in\mathcal E_{\chi_!}$.

On a $S^n=S^n_+\cup S^n_-$ et $S^n_+\cap S^n_-\simeq S^{n-1}$, \inserted{Or $S^n_+\simeq S^n_-\simeq B^n$ donc $\in\mathcal E_\chi$} donc on conclut de c) par réc. sur $n$ que $S^n\in\mathcal E_\chi$ (\uline{NB} $S^0=\{\mathrm{pt}\}\sqcup\{\mathrm{pt}\}\in\mathcal E_\chi$ par d) et e), et $\chi(S^0)=1+1=2$), et on trouve $\chi(S^n)=\chi(S^n_+)+\chi(S^n_-)-\chi(S^{n-1})\deleted{\unsure{+}}=2-\chi(S^{n-1})$, d'où par récurrence sur $n$ $\chi(S^n)=1+(-1)^n$. Comme $S^n$ cpt, $S^n\in\mathcal E_{\chi_!}$ et $\chi_!(S^n)=\chi(S^n)$ par (a). Enfin, $R^n\simeq S^n-\mathrm{pt}$, donc par (b) $R^n\in\mathcal E_{\chi_!}$ et $\chi(R^n)=\chi(S^n)-\chi(\mathrm{pt})=(1+(-1)^n)-1=(-1)^n$, cqfd.

\textcircled{3} Soit $X_{-1}=\varnothing\subset X_0\subset X_1\subset\cdots\subset X_n=X$ filtration de l'espace $X$, tel que $\forall 0\leq i\leq n$, $X_i-X_{i-1}$ \deleted{\unsure{est} \unsure{$\mathbb R$}} a un nb fini $f_i$ de comp. connexes ttes homéom. à $\mathbb R^i$ (i.e. $X_i-X_{i-1}\simeq f_i\times\mathbb R^i$). Alors $X\in\mathcal E_{\chi_!}$, et on a :
\[
\chi_!(X)=\sum(-1)^i f_i
\]
(donc si $X$ est cpt, on a $\deleted{\unsure{\chi}}X\in\mathcal E_\chi$ et $\chi(X)=\sum(-1)^i f_i\ \ldots$

% archive sheet number: 57

\scanpage{59}
\mspage{9.}

4.3. - \uline{La politique de recherche des établissements}

Un établissement d'enseignement supérieur possède en général un Conseil scientifique. Il lui revient de définir les grandes lignes de la politique scientifique de l'établissement. L'objectif est clair :

- compte tenu des moyens en bâtiments, personnel enseignant et technicien, comment utiliser les crédits spécifiques pour la recherche de façon optimum, pour faire la meilleure recherche, cel\deleted{\unsure{a|le}}le qui assure le renon de l'établissement et lui permet de s'associer avec les grands établissements de recherche, en premier lieu avec le C.N.R.S.

- Il appartient au Conseil scientifique de conduire la politique de recherche de l'établissement, sous la responsabilité du Président. La recherche est l'affaire de l'ETABLISSEMENT, chacun étant solidaire pour la mener à bien.

Une politique recherche est, typiquement, celle que fait naitre des équipes qui deviennent peu à peu des équipes associées avec le C.N.R.S. et avec l'I.N.S.E.R.M.. La collaboration prolongée de l'établissement à frais partagés avec un organisme de recherche est le critère de réussite d'une politique de recherche.

5 - \uline{La politique "recherche" de la D.E.S.U.R.}

5.1. \uline{La D.E.S.U.R. doit remplir une mission de complexe dans le cadre du plan-recherche lui permettant} :

a) d'agir "en aval" au niveau des établissements en respectant leur autonomie ; elle ne peut donner des ordres ni définir des programmes \uline{a priori} ; elle doit par contre, effectuer un contrôle a postériori, pour avoir des éléments de jugement sur l'activité de recherche des établissements.

b) d'agir "en amont" au niveau national en tenant sa place, parmi une dizaine d'organismes fortement hièrarchisés, pour obtenir sa part de l'enveloppe-recherche. La D.E.S.U.R. doit être "crédible" par la D.G.R.S.T.

% archive sheet number: 58

\scanpage{60}

\textcircled{f}
\[
\begin{aligned}
X,Y\in\unsure{\mathcal E|\mathscr E}_{\chi_!}&\Rightarrow X\times Y\in\mathcal E_{\chi_!} &&\text{et on a} &\chi_!(X\times Y)&=\chi_!(X)\chi_!(Y)\\
X,Y\in\mathcal E_\chi&\Rightarrow X\times Y\in\mathcal E_\chi &&\text{et on a} &\chi(X\times Y)&=\chi(X)\chi(Y).
\end{aligned}
\]

(La première assertion résulte de Künneth pour $H^*_!(X\times Y,k)$, la deuxième de Künneth pour $H_*(X\times Y,k)$, dont le dual est $H^*(X\times Y,k)$, donc qui permet de tester l'appartenance à $\mathcal E_\chi$ et la valeur de $\chi$ $\ldots$)

\textcircled{g} Soit $f:X\unsure{-|\longrightarrow}Y$ fibration loc. triviale de fibre $F$.

Si $Y,F\in\mathcal E_{\chi_!}$ ; alors $X\in\mathcal E_{\chi_!}$ et on a
\[
\chi_!(X)=\chi_!(Y)\chi_!(F).
\]
\margin{avec léger grain de sel}% slanted

Énoncé analogue pour \inserted{$Y,F\in$} $\mathcal E_\chi$ et \deleted{\unsure{$\chi(X)=$}}
\[
\chi(X)=\chi(Y)\chi(F)
\]

\uline{Dém.} On a
\[
H^*_!(X)\Leftarrow E_2^{pq}=H^p_!\bigl(Y,\underbrace{R^q f_!(k_X)}_{\substack{\text{syst. local sur }Y\\\text{à fibres }\sim H^q_!(F)}}\bigr)
\]
donc \uline{sous réserve} qu'on sache que $E_2^{**}$ de dim finie, on trouve $\deleted{\unsure{\chi}}X\in\mathcal E_\chi$ et
\[
\begin{aligned}
\chi(X)&=\sum_{pq}(-1)^{p+q}\operatorname{rg}_k H^p_!(Y,R^q f_!(k_X))\\
&=\sum_q(-1)^q\sum_p\underbrace{(-1)^p\operatorname{rg}_k H^p_!(Y,R^q f_!(k_X))}_{b^q_!(F)\chi_!(Y)}\\
&=\chi_!(Y)\deleted{\unsure{\chi}}\sum_q(\deleted{\unsure{+}}-1)^q b^q_!(F)=\chi_!(Y)\chi_!(F)
\end{aligned}
\]

Il faut prouver

\uline{Lemme} Soit $\mathcal R$ syst. local de $k$-vectoriels sur \deleted{\unsure{$Y$|$X$} \unsure{loc. cpt}}, $Y\in\mathcal E_{\chi_!}$, à fibres de rang \inserted{fini} $b$. Alors
\[
\chi_!(Y,\mathcal R)\overset{\mathrm{déf}}{=}\sum_p(-1)^p H^p(Y,\mathcal R)
\]
existe (i.e. $\unsure{\sum|\bigoplus}_p H^p(Y,\mathcal R)$ de dim finie sur $k$) et on a
\[
\chi_!(Y,\mathcal R)=b\chi_!(Y)
\]

Marche tout au moins si $Y$ \deleted{triangulable} \inserted{admet filtration finie} avec $Y_i-Y_{i-1}$ somme finie de cellules \inserted{$Z_{ij}$} $\simeq R^{\unsure{d|\delta}_{ij}}$, car
\margin{\textcircled{\unsure{gr.|gn|gr}}}
On aura
\[
\begin{aligned}
\chi_!(Y,\mathcal R)&=\sum_{ij}\chi_!(Z_{ij},\underset{(\text{ct. sur }Z_{ij})}{\mathcal R})=\sum_{ij}b\chi_!(Z_{ij})\\
&=b\sum\chi_!(Z_{ij})\unsure{\mathpunct{:}|=}b\chi_!(Y).
\end{aligned}
\]

% archive sheet number: 59

